\documentclass[10pt,a4paper]{article}
\usepackage{amsmath,amssymb,amsthm,graphicx}
\RequirePackage{amsthm,amssymb}
\usepackage{amsmath}
\usepackage[hidelinks]{hyperref}
\newcommand{\om}{\omega}
\newcommand{\oequiv}{\approx}
\newcommand{\setconcat}{+}
\newcommand{\COL}{\operatorname{COL}}
\newcommand{\N}{\mathbb{N}}
\newcommand{\Z}{\mathbb{Z}}

\newcommand{\todo}[1]{\textbf{TODO:} #1}

\newcommand{\capcrules}{Coloring Rules}
\newcommand{\crulenospace}{CR}
\newcommand{\cruleLong}{coloring rule}
\newcommand{\crulesLong}{coloring rules}
\newcommand{\crule}{\crulenospace~}
\newcommand{\crules}{\crulenospace s~}
\newcommand{\crulesnospace}{\crulenospace s}

\newcommand{\capgencrules}{General Coloring Rules}
\newcommand{\gencrulenospace}{GCR}
\newcommand{\gencruleLong}{general coloring rule}
\newcommand{\gencrulesLong}{general coloring rules}
\newcommand{\gencrule}{\gencrulenospace~}
\newcommand{\gencrules}{\gencrulenospace s~}
\newcommand{\gencrulesnospace}{\gencrulenospace s}

\newcommand{\xleq}{\preceq_{\overline{A}}}
% \newcommand{\assny}{\mathcal{Y}}
\newcommand{\assny}{\overline{B}}
\newcommand{\assnc}{\overline{C}}

\newcommand{\case}[1]{\textbf{Case #1.}}

\newcommand{\numpoints}{n}
\newcommand{\inumpoints}{i}
\newcommand{\idegr}{j}
\newcommand{\numcolors}{c}
\newcommand{\commonconst}{a}
\newcommand{\Commonconst}{A}
\newcommand{\rareconst}{b}
\newcommand{\Rareconst}{B}

\theoremstyle{plain}
\newtheorem{theorem}{Theorem}[section]
\newtheorem{lemma}[theorem]{Lemma}

\theoremstyle{definition}
\newtheorem{definition}[theorem]{Definition}
\title{Exact big Ramsey degree four for two-element subsets of the ordered set of integers}
\author{Joanna Boyland, William Gasarch, Nathan Hurtig, Robert Rust}
\date{arxiv:2305.07192v3; CC BY 4.0}
\begin{document}
\maketitle
\noindent\textit{Editorial note.} The counted statement and its proof are adjacent in the source: the unlabelled proof environment beginning at source line 498 is the proof of Theorem the:Z2 printed immediately above it, and the article gives no other proof of this theorem.\par\medskip
\noindent\textit{Editorial note.} The retained context includes the definition of the big Ramsey degree, which ends with a citation of Kechris, Pestov and Todorcevic at source line 227. The bibliography fragment retained for this source supplies that single entry together with the two further entries of the article, so the citation resolves and no reference is left undefined.\par\medskip
\noindent\textit{Editorial note.} The original numbers theorems by section. The assembled wrapper keeps the author's own theorem environment, so the retained section opening from source line 465 keeps the numbering context of the original without inventing a new number.\par\medskip
\noindent\textit{Editorial note.} Only the counted theorem is a counting unit. The n = 1 sibling and the general integers theorem T(n,zeta) = 2\textasciicircum{}n are not counted: the general theorem's printed lower bound is explicitly left to the reader at source lines 588-589, while the counted n = 2 theorem has a complete printed proof of both bounds.\par\medskip
\noindent\textit{Editorial scope.} One counted exact value, the big Ramsey degree T(2,zeta) = 4 for two-element subsets of the ordered set of integers, with both bounds of the article's complete adjacent proof. The n = 1 sibling, the article's general integers theorem, the omega-squared example, the omega\textasciicircum{}d and omega\textasciicircum{}d times k theorems and the general theorem for ordinals below omega\textasciicircum{}omega are not counted.\par\medskip
\noindent\textit{Allowed standard prerequisites (editorial).}\begin{itemize}
\item Infinite Ramsey theorem, in the article's own big-Ramsey-degree form. For every n and every finite colouring of the n-element subsets of the naturals there is an infinite subset on which the colouring is constant; the article restates this as T(n, omega) = 1 at source lines 241-251.. source lines 239-251, retained as context. Applied twice in the counted proof, to the four-tuple colouring of the pairs of omega (source line 507) and, in the sibling arguments of the same section, to colourings of the naturals.
\item Order equivalence and copies of the integers. Two ordered sets are order-equivalent when an order-preserving bijection exists between them; the ordered set zeta is the integers with their natural order, and a subset of zeta is order-equivalent to zeta exactly when it has neither a least nor a greatest element and is unbounded in both directions.. definitions at source lines 172-196 and 213-216, retained as context. Both directions of the counted proof rebuild or test a subset H of zeta that is order-equivalent to zeta, and the lower bound uses that such an H must contain positive elements, negative elements and pairs of both signs.
\item Elementary manipulations of Z with signs, absolute values and parity. Every integer is either a nonnegative integer or the negative of one; the absolute value of a negative integer is its negation; an infinite subset of the naturals can be split by the parity of its indices.. source lines 501-513 and 515-537. Building the four sign combinations of a pair, comparing absolute values in the four cases of the lower-bound colouring and alternating the signs of the sequence h\_0, h\_1, ... when the copy of zeta is rebuilt.
\end{itemize}\medskip
\begin{definition}
Let $\N$ be the set of natural numbers, including 0. 
Let $\numpoints,\numcolors \in \N$ and $S$ be a set. Then we define $[\numcolors]$ to be $ \{ 0, \ldots, \numcolors-1\}$
(note $0 \in [\numcolors]$, but $\numcolors \not \in [\numcolors]$). Note that if $\numcolors=0$ then $[\numcolors] = \emptyset$. 
We use $\binom{S}{\numpoints}$ to denote the set of all $\numpoints$-element subsets of $S$. 
\end{definition}

\begin{definition}
    A function $\COL\colon \binom{S}{\numpoints} \to [\numcolors]$ is called a \textit{coloring}.
    A set $H\subseteq S$ is called \textit{homogeneous} if $\COL$ restricted to 
    $\binom{H}{\numpoints}$ is constant,
    that is, $\left|\COL\left(\binom{H}{\numpoints}\right)\right|\le 1$.
\end{definition}

\begin{definition} 
Let $\mathcal{A} = (A,\preceq_A)$ and $\mathcal{B} = (B,\preceq_B)$ be ordered sets, where
$\preceq_A, \preceq_B$ are total orders on $A, B$ respectively.

\begin{enumerate}

\item Two ordered sets
$\mathcal{A},\mathcal{B}$ are \textit{order-equivalent}, denoted $\mathcal{A}\oequiv \mathcal{B}$, 
if there exists an order-preserving bijection $f\colon A \to B$ such that, for all $a_1, a_2 \in A$: 

$$a_1 \preceq_A a_2 \iff f(a_1) \preceq_B f(a_2).$$

\item 
The ordered set $\mathcal{A} \setconcat \mathcal{B}$ is defined to be
$(A\sqcup B, \preceq),$ where the disjoint union $A\sqcup B$ is formally defined to be
$(\{0\} \times A) \, \cup\, (\{1\} \times B)$, and where $\preceq$ is the dictionary ordering on $A\sqcup B.$
Note that $\setconcat$ agrees with the definition of ordinal addition, and is not commutative in general.

\item 
Let $\mathcal{A}=(A,\preceq_A)$ be an ordered set with $A \subseteq \Z$. Then define $-\mathcal{A}$ to be $(-A,\preceq_{-A})$, where 
$-A = \{ -a \colon a \in A\}$, and
$-a \preceq_{-A} -b \iff b \preceq_{A} a$.
\end{enumerate} 

\end{definition} 

\begin{definition}\label{def:homogeneous}
Let $S$ be an ordered set, $S'\subseteq S$, and $\numpoints,\numcolors,t \in \N.$ Let $\COL \colon \binom{S}{\numpoints} \to [\numcolors]$ be a coloring.
\begin{enumerate}
    \item 
$S'$ is \textit{homogeneous} if  $\left|\COL\left(\binom{S'}{\numpoints}\right)\right|= 1$. 
\item 
 $S'$ is \textit{$t$-homogeneous} if $\left|\COL\left(\binom{S'}{\numpoints}\right)\right|\leq t$. 
 \item 
$S'$ is $S$-\textit{$t$-homogeneous} if both $S'$ is $t$-homogeneous and $S'\oequiv S$.
\end{enumerate}
\end{definition}

\begin{definition}
We identify every ordinal $\alpha$ as the ordered set of all ordinals $\beta < \alpha$, with 0 as the least ordinal. Additionally, let $\omega$ be the ordered set containing only the naturals, $\zeta$ be the ordered set containing only the integers, and $\eta$ be the ordered set containing only the rationals, all under their respective natural orderings.

\end{definition}

\begin{definition}\label{def:T}
Let $S$ be an ordered set.
For any $\numpoints\in \N$, $T(\numpoints,S)$
is the least $t\in \N$ such that for all constants $\numcolors \geq 1$ and
colorings  $\COL \colon \binom{S}{\numpoints} \to [\numcolors]$, there exists some $S' \subseteq S$ 
such that $S'$ is $S$-$t$-homogeneous.
Note that $t$ is independent of $\numcolors$.
If no such $t$ exists then we define $T(\numpoints,S)=\infty$. 
The quantity $T(\numpoints,S)$ is called the \textit{big Ramsey degree} of $\numpoints$ in $S$.
The term was first coined by Kechris et al.~\cite{KPT-2005}. Other sources use the Ramsey arrow notation modified for ordered sets. In that notation, $T(\numpoints,S)=t$ is short-hand for both
$S \to (S)^{\numpoints}_{\numcolors,t}$ and $S \not \to
(S)^{\numpoints}_{T(\numpoints,S),t-1}$ for all $\numcolors \in \N.$
We use the $T(n, S)$ notation in this paper.
\end{definition}

Ramsey's Theorem on $\N$ gives an infinite 1-homogeneous subset of $\N$. Theorem~\ref{th:ramsey} restates Ramsey's Theorem using big Ramsey degree notation.

\begin{theorem}\label{th:ramsey}
For all $\numpoints \in \N,$
% \begin{enumerate}
    % \item 
$T(\numpoints,\omega)=1.$
% \item 
% Let $\numcolors \geq 1$ and $\COL\colon\binom{\omega}{\numpoints}\to [\numcolors]$.
% Then there exists some $H\oequiv \omega$ such that 
% $$\left | \COL\left(\binom{H}{\numpoints}\right) \right | = 1.$$
% \end{enumerate}
\end{theorem}

\section{Big Ramsey Degrees of \texorpdfstring{$\zeta$}{Z}}\label{se:Z}


Our general result is $T(n, \zeta) = 2^n$.
We first prove $T(1,\zeta)=2$
and $T(2,\zeta)=4$. Both of these will 
build notation and structure toward the main result.

\begin{theorem}\label{the:Z2}
The big Ramsey degree $T(2,\zeta)=4$. 
\end{theorem}

\begin{proof}
We first prove $T(2,\zeta)\leq 4$. Let $\numcolors \in \N$ and let $$\COL\colon \binom{\zeta}{2}\to [\numcolors]$$
be a $\numcolors$-coloring of $\binom{\zeta}{2}$.
Let $\COL'\colon \binom{\omega}{2} \to [\numcolors]^4$ be defined by 

$$
\COL'(x,y) = (\COL(-x,-y), \COL(-x,y), \COL(x,-y), \COL(x,y))
$$ where $x < y.$

By Theorem~\ref{th:ramsey} there exists an $\om$-1-homogeneous set $G$. 
Let the color of the homogeneous set be $(c_1,c_2,c_3,c_4)$.
Let $G = \{h_0 < h_1 < \cdots \}$. 

Then consider the set $$H= \{-h_i \colon i \textnormal{ is even}\} \setconcat \{h_i \colon i \textnormal{ is odd}\},$$ which is order-equivalent to $\zeta$. Let $s_in_i, s_jn_j \in H,$ where $n_i,n_j\geq 0$ and $s_i,s_j \in \{-1,1\}$. Then by definition of $\COL'$, $\COL(s_in_i,s_jn_j)$ is in $\{c_1, c_2, c_3, c_4\}$, depending on $s_i$ and $s_j$. As an aside, this method only works when $n_i$ and $n_j$ are guaranteed to be distinct, which we forced with our alternation of sign by the parity of index.
Since $\COL(H)\subseteq \{c_1,c_2,c_3,c_4\}$, $H$ is $\zeta$-4-homogeneous.
Therefore $T(2,\zeta)\leq 4$.

We now prove $T(2,\zeta)\geq 4$. Let $\COL\colon\binom{\zeta}{2}\to [4]$ be the coloring
\begin{align*}
\COL(x,y)=
\begin{cases}
0 & \text{$x,y\geq 0$} \cr
1 & \text{$x\geq 0$, $y<0$, and $|x|\leq |y|$} \cr
2 & \text{$x\geq 0$, $y<0$, and $|x|> |y|$} \cr
3 & \text{$x<0$, $y<0$.} \cr
\end{cases}
\end{align*}
Let $H$ be an arbitrary order-equivalent subset of $\zeta.$
The following $x,y$ are guaranteed to exist because $H \subseteq \zeta$ and $H \oequiv \zeta$:
\begin{enumerate} % T(2,\zeta) lower bound
\item
Let $x,y\in H$ such that $x,y\ge 0$. Then $\COL(x,y)=0$. 
\item
Let $x,y\in H$ such that $x\ge 0$, $y<0$, and $|x|\le |y|$. Then $\COL(x,y)=1$. 
\item 
Let $x,y\in H$ such that $x\ge 0$, $y<0$, and $|x|>|y|$. Then $\COL(x,y)=2$. 
\item
Let $x,y\in H$ such that $x<0$, $y<0$. Then $\COL(x,y)=3$. 
\end{enumerate}  % T(2,zeta) lower bound 
Therefore $T(2,\zeta)\geq 4$, and with the previous result, $T(2,\zeta) = 4$.
% We show there is no $\zeta$-3-homogeneous set. Let $H \subseteq \zeta$ with $H \oequiv \zeta.$ Then there is some order-preserving bijection between the two sets. Let $h_0$ be the element of $H$ that corresponds to $0 \in \zeta.$ Because $H \subseteq \zeta,$ we can compare elements of $H$ to $0 \in \zeta$ and also calculate their absolute values in $\zeta.$ Take successors of $h_0$ until some $h_+ \in H$ with $h_+ > 0$ is found. Likewise, find some $h_- \in H$ with $h_- < 0$ using predecessors. Take successors of $h_+$ until finding some $h_{++}$ with $|h_{++}| > |h_{-}|$ and take predecessors of $h_-$ until finding some $h_{--}$ with $|h_{--}| \geq |h_+|.$ Then \begin{itemize}
%     \item $\COL(h_{++},h_+) = 1,$
%     \item $\COL(h_{+},h_{--}) = 2,$
%     \item $\COL(h_{++},h_{-}) = 3,$ and
%     \item $\COL(h_{-},h_{--}) = 4$ by the definition of $\COL.$
% \end{itemize} Therefore $\COL$ outputs 4 colors on every $H \subseteq \zeta$ with $H \oequiv \zeta$ so no $\zeta$-3-homogeneous set exists. Therefore $T(2,\zeta)\geq 4$, and with the previous result, $T(2,\zeta) = 4$.
\end{proof}

\begin{thebibliography}{9}
\bibitem{BDR21}
A.~Blass, N.~Dobrinen, and D.~Raghavan.
\newblock The next best thing to a {P}-point.
\newblock {\em The Journal of Symbolic Logic}, 80(3):866–900, 2015.
\bibitem{KPT-2005}
A.~S. Kechris, V.~G. Pestov, and S.~Todorcevic.
\newblock Fra{\"i}ss{\'e} limits, {R}amsey theory, and topological dynamics of
  automorphism groups.
\newblock {\em Geometric {\&} Functional Analysis GAFA}, 15(1):106--189, Feb
  2005.
\bibitem{Ramsey-1930}
F.~Ramsey.
\newblock On a problem of formal logic.
\newblock {\em Proceedings of the London Math Society}, 30(1):264--286, 1930.
\end{thebibliography}
\end{document}
